\section{What the finite-cokernel construction still requires uniformly}
\label{sec:uniform-positive-gate}

This section addresses item 4 of the continuation prompt.  It proves a
size obstruction to the literal implementation of G-040 and identifies
the remaining representation problem.  It does not prove that a
polynomial-size positive formula is impossible.

Let $m\ge3$, let $W=\mathfrak h_m$ be the sum-zero subrepresentation
of $\mathbb Q^m$, and put $P=\operatorname{Sym}W$.  The invariant
form identifies $W$ with $W^*$ throughout.  Give the two tensor
factors of
\[
 \phi_b:P\otimes\operatorname{Sym}^bW
       \longrightarrow P\otimes\operatorname{Sym}^{b-1}W,
 \qquad
 f\otimes v_1\cdots v_b\longmapsto
 \sum_{i=1}^b fv_i\otimes\prod_{j\ne i}v_j
 \tag{PNEW-G1}
\]
their total polynomial degree, and let $C=\operatorname{coker}\phi_b$.
For $\nu\vdash m$ put $C_\nu=\operatorname{Hom}_{S_m}([\nu],C)$.
Write $e_i$ for the degree-$i$ elementary symmetric polynomial
on $W$, and set
\[
 B=\min(m,b-1),\qquad J=\mathbb Q[e_2,\ldots,e_B].
 \tag{PNEW-G2}
\]
These are the actual modules and variables in G-040, not auxiliary
modules with the same Hilbert functions.

\begin{proposition}[PNEW-010: a superpolynomial minimal cover]
For every integer $r\ge2$, take
\[
 m=2r,\qquad b=r(r+1)+1,\qquad \nu=(m).
 \tag{PNEW-G3}
\]
Every homogeneous generating set of the $J$-module $C_\nu$ has
at least $2^{r-1}$ members.  In particular the minimal cover
$F_0\longrightarrow C_\nu$ in G-040 has rank at least $2^{r-1}$.
Its rank is not bounded by any polynomial in $m+b$.
\end{proposition}
\begin{proof}
The target of (PNEW-G1) has no component below degree $b-1$;
in degree $b-1$ it is exactly $\operatorname{Sym}^{b-1}W$.
The domain has no component at that degree.  Consequently
\[
 C_{b-1}=\operatorname{Sym}^{b-1}W,
 \qquad C_d=0\quad(d<b-1).
 \tag{PNEW-G4}
\]
Since every variable of $J$ has positive degree, the entire
degree-$b-1$ component of $C_\nu$ injects into
$C_\nu/J_+C_\nu$.  The images of any generating set must span
that quotient, and hence must span this component.

The fundamental theorem of symmetric polynomials gives
\[
 (\operatorname{Sym}W)^{S_m}=\mathbb Q[e_2,\ldots,e_m].
 \tag{PNEW-G5}
\]
Here is the passage to $W$ explicitly.  In the polynomial ring
on $\mathbb Q^m$, the equation of $W$ is $e_1=0$.
An invariant residue class has an invariant lift by averaging
over $S_m$.  An invariant multiple $e_1g$ has an invariant
quotient $g$, since $e_1$ is invariant and the polynomial ring
is a domain.  Thus the invariant quotient is exactly
$\mathbb Q[e_1,\ldots,e_m]/(e_1)$, proving (PNEW-G5), including
algebraic independence of the remaining generators.

Put $N=r(r+1)=b-1$.  For every subset
$A\subseteq\{2,\ldots,r\}$ consider the invariant monomial
\[
 M_A=e_2^{\left(N-\sum_{j\in A}2j\right)/2}
                 \prod_{j\in A}e_{2j}.
 \tag{PNEW-G6}
\]
Its exponent is an integer because $N$ is even.  It is
nonnegative because
$\sum_{j=2}^r2j=r(r+1)-2<N$.  Every index $2j$ is at most
$m=2r$, and every $M_A$ has weighted degree $N$.
Distinct subsets give distinct monomials in the algebraically
independent generators of (PNEW-G5).  The $2^{r-1}$ such
monomials are therefore linearly independent vectors of
$(C_{(m)})_N$ and of $C_{(m)}/J_+C_{(m)}$.
This proves the generating-set bound.  A minimal homogeneous
cover has one free generator for each homogeneous basis vector
of the latter quotient, so it obeys the same bound.
Finally $m+b=r^2+3r+1$, whereas $2^{r-1}$ exceeds every fixed
power of that quantity as $r\to\infty$.
\end{proof}

\paragraph{The exact computational consequence.}
The obstruction occurs before G-040's Hermann or Gr\"obner degree
estimate.  A procedure which explicitly lists these minimal
generators, or forms the corresponding full matrices, cannot
have construction size polynomial in varying $m,b$ merely by
improving those degree estimates.  This is a lower bound on that
specified representation of the algorithm, not on the
plethysm function.  A positive $\#P$ formula may have
exponentially many accepting witnesses, each with a short
description; the subsets in (PNEW-G6) themselves illustrate
precisely that distinction.

\begin{proposition}[The finite Schur pieces return a stable plethysm family]
For $0\le j\le b-1$ one has
\[
 C_{b-1+j}\simeq S_{(b-1,j)}W.
 \tag{PNEW-G7}
\]
Set $n=b-1$ and
\[
 \lambda=((m-1)(b-1)-j,\ b-1,\ j,\ 0^{m-3}).
 \tag{PNEW-G8}
\]
Then $\lambda\vdash mn$ is a partition and, with the ordinary
outer-inner-output convention,
\[
 \dim\operatorname{Hom}_{S_m}([\nu],C_{b-1+j})
           =p_{\nu,(n)}^{\lambda}.
 \tag{PNEW-G9}
\]
\end{proposition}
\begin{proof}
At total degree $b-1+j$, the target of (PNEW-G1) is
$\operatorname{Sym}^{j}W\otimes\operatorname{Sym}^{b-1}W$.
For $j=0$ the domain is zero.  For $j\ge1$ the domain is
$\operatorname{Sym}^{j-1}W\otimes\operatorname{Sym}^{b}W$.
The two-factor Pieri rule gives each with multiplicity one.
Using two pairs of variables, the highest vector in the domain's
summand of second part $i$ is
\[
 x^{j-1-i}X^{b-i}(xY-yX)^i,
 \qquad 0\le i\le j-1.
\]
The operator $x\partial_X+y\partial_Y$ annihilates the determinant
$xY-yX$ and maps this vector to $b-i$ times the target's
highest vector.  Since $b-i>0$, each common summand is mapped
isomorphically.  The only remaining target summand has second
part $j$ and first part $b-1$, proving (PNEW-G7).

The first gap in (PNEW-G8) is
$(m-2)(b-1)-j\ge0$; the second is $b-1-j\ge0$.
The sum is $m(b-1)$.  In the rechecked actual balanced-space
map G-015, equations (GF.4)--(GF.7), the tail is
$(b-1,j,0,\ldots,0)$ and the point-ideal exponent is
$\max(0,\lambda_2-n)=0$.  Hence that map identifies
$(S_\lambda\mathbb Q^m)_{(n^m)}$ with $S_{(b-1,j)}W$,
with its literal $S_m$ action.  Taking the $[\nu]$-covariant
space proves (PNEW-G9).  The balanced-space identification is
an inherited theorem whose whole required proof was read here;
the present argument does not silently replace it by equality
of dimensions.
\end{proof}

Thus the semistandard-tableau basis of the Schur module in
(PNEW-G7) is not yet a rule for its $S_m$ multiplicities.
Its covariant multiplicities already form the stable plethysm
subfamily (PNEW-G9).  This is an exact reduction, not a
complexity-hardness result.

\begin{lemma}[A permutation-basis shortcut already fails at $b=2$]
There is no basis of $C_1=W$ that is permuted by every element
of $S_m$.
\end{lemma}
\begin{proof}
An $m$-cycle has trace zero on $\mathbb Q^m$ and trace one on
its constant line, hence trace $-1$ on $W$.
The trace of a permutation of a basis is its nonnegative
number of fixed basis vectors.  This is a contradiction.
\end{proof}

This lemma excludes only a literal permutation basis.  It does
not exclude a basis compatible with Specht multiplicities,
a signed basis, or a positive counting rule of another kind.
Similarly a basis of $C$ as a vector space does not by itself
give the $J$-covariants, their minimal relations, or the kernel
$K_J$ in G-040.  These are distinct operations in the actual
construction.

\paragraph{Literature comparison and remaining gate.}
Orellana and Zabrocki's multiset-tableau rule, Theorem 3.1 of
\emph{A combinatorial model for the decomposition of multivariate
polynomial rings as $S_n$-modules}, applies to tensor products
of symmetric and exterior powers of the permutation
representation.  Its theorem and the full displayed proof were
read.  It counts the pertinent multidegree modules; it does not
give a positive rule for the kernel of the transfer above.
Narayanan, Paul, Prasad, and Srivastava's Littlewood formula
(2021, Theorem 2.6) identifies the ordinary permutation-module
restriction coefficient with
$\langle s_\alpha,s_\nu[1+h_1+h_2+\cdots]\rangle$.
That identity likewise does not itself give a positive rule for
the required Schur restriction.  Lee's 2026 result supplies a
combinatorial rule for shapes with at most three columns
(Theorem 1.3); shapes $(b-1,j)$ have unboundedly many columns
here.  These exact source scopes have been checked, and no
claim that the entire current literature rules out another
construction is intended.

The unresolved gate is a uniformly short, polynomial-time
verifiable indexing of the actual covariant and relation
contributions, compatible with their multiplication maps, or a
different direct positive rule for the original stable
two-row-tail multiplicities.  No such indexing, nor the
requested three-row-tail extension, is supplied in this section.

\paragraph{Source and verification locators.}
The inherited algorithm is
\path{sources/RETURN_GCT_PLETH_checkpoint_09_20261121T0546Z/work/root/fixed_second_part_report.md},
G-040, especially (FS.3)--(FS.7) and \S\S2--6.
The actual transport is the adjacent
\path{general_fiber_model.md}, G-015, (GF.1)--(GF.13).
Both complete required proofs were read.  The displayed new
lower bound is a symbolic proof for all $r$, not a finite
experimental extrapolation.  The companion script
\path{scripts/check_uniform_positive_gate.py} gives independent
finite calibration values for the generator bound and the
Cox grading dictionary; these computations are supplementary.

\paragraph{References.}
Lee, M. (2026). Restriction coefficients for partitions with at
most three columns. \emph{Algebraic Combinatorics, 9}(2),
543--555. \url{https://doi.org/10.5802/alco.476}.
Primary proof version: \url{https://arxiv.org/pdf/2506.11837v2}.

Narayanan, S. P., Paul, D., Prasad, A., \& Srivastava, S.
(2021). Polynomial induction and the restriction problem.
\emph{Indian Journal of Pure and Applied Mathematics, 52},
643--651. \url{https://doi.org/10.1007/s13226-021-00185-7}.

Orellana, R., \& Zabrocki, M. (2020).
\emph{A combinatorial model for the decomposition of multivariate
polynomial rings as $S_n$-modules} (arXiv:1906.01125v2).
\url{https://arxiv.org/pdf/1906.01125v2}.
